\section{Sistemas híbridos y enfoques basados en recuperación de información}
\textcolor{teal}{Descripción de aproximaciones que combinan modelos lingüísticos, técnicas tradicionales y recuperación de información. Revisión de los artículos académicos orientados al uso de sistemas de evaluación co nrecuperación deinformación y recomendaciones.}

\textcolor{teal}{
    2.4.1 Sistemas híbridos
    Qué son:
    \begin{itemize}
        \item combinación de:
        \begin{itemize}
            \item reglas lingüísticas
            \item modelos clásicos
            \item LLMs
        \end{itemize}        
    \end{itemize}
    Ventajas:
    \begin{itemize}
        \item mayor control
        \item más interpretabilidad
        \item menor dependencia del modelo
    \end{itemize}    
    Ejemplo aplicado: corrector gramatical + LLM evaluador
}

\textcolor{teal}{
    2.4.2 RAG (Retrieval-Augmented Generation)
    Qué explicar:
    \begin{itemize}
        \item uso de información externa
        \item mejora de contexto
    \end{itemize}
    Aplicación en evaluación:
    \begin{itemize}
        \item enunciado
        \item solución modelo
        \item rúbrica
    \end{itemize}
    Ventajas:
    \begin{itemize}
        \item reduce alucinaciones
        \item mejora precisión
    \end{itemize}
    Problemas:
    \begin{itemize}
        \item complejidad
        \item dependencia de datos externos
    \end{itemize}
}

\textcolor{teal}{
    2.4.3 LLMs como evaluadores (LLM-as-a-judge)
    Qué explicar:
    \begin{itemize}
        \item uso de LLMs para evaluar texto
    \end{itemize}
    Problemas identificados en literatura:
    \begin{itemize}
        \item sesgo
        \item inconsistencia
        \item sensibilidad al prompt
    \end{itemize}
    Conexión con tu TFM -> Tu trabajo analiza precisamente estos problemas
}